\documentclass[12pt,a4paper]{article}

\input{glyphtounicode}
\pdfgentounicode=1
\usepackage{cmap}
\usepackage[utf8]{inputenc}
\usepackage[T1]{fontenc}
\usepackage{lmodern}
\usepackage{amsmath,amssymb,amsthm,mathtools}
\usepackage[colorlinks=true,linkcolor=blue,citecolor=red,urlcolor=blue]{hyperref}
\hypersetup{
  pdftitle={The Spectral Zookeeper: A Conditional Reduction toward the Riemann Hypothesis via Spectral Microcluster Closure in the Connes--Consani--Moscovici Framework},
  pdfauthor={Lukas Geiger},
  pdfsubject={Preprint on a conditional reduction toward the Riemann Hypothesis via spectral microcluster closure},
  pdfkeywords={Riemann Hypothesis, Weil explicit formula, noncommutative geometry, spectral approximation, Abel summation, Cauchy interlacing}
}
\usepackage{geometry}
\geometry{margin=2.5cm}
\usepackage{enumitem}
\usepackage{booktabs}
\usepackage{graphicx}
\usepackage[capitalise,noabbrev]{cleveref}
\emergencystretch=2em

\newtheorem{theorem}{Theorem}[section]
\newtheorem{lemma}[theorem]{Lemma}
\newtheorem{proposition}[theorem]{Proposition}
\newtheorem{corollary}[theorem]{Corollary}
\newtheorem{conjecture}[theorem]{Conjecture}
\newtheorem{hypothesis}[theorem]{Hypothesis}
\theoremstyle{definition}
\newtheorem{definition}[theorem]{Definition}
\theoremstyle{remark}
\newtheorem{remark}[theorem]{Remark}

% === Custom commands ===
\newcommand{\CC}{\mathbb{C}}
\newcommand{\RR}{\mathbb{R}}
\newcommand{\ZZ}{\mathbb{Z}}
\newcommand{\NN}{\mathbb{N}}
\newcommand{\HH}{\mathcal{H}}
\newcommand{\PP}{\mathcal{P}}
\newcommand{\Cl}{\mathrm{Cl}}
\newcommand{\reg}{\mathrm{reg}}
\newcommand{\off}{\mathrm{off}}
\newcommand{\rank}{\mathrm{rank}}
\newcommand{\sgn}{\mathrm{sgn}}
\newcommand{\Tr}{\mathrm{Tr}}
\newcommand{\ip}[2]{\langle #1, #2 \rangle}
\newcommand{\norm}[1]{\| #1 \|}
\newcommand{\abs}[1]{| #1 |}
\newcommand{\utilde}{\tilde{u}}
\newcommand{\Ctilde}{\tilde{C}}
\newcommand{\wmin}{w_{\min}}

\title{The Spectral Zookeeper:\\[4pt]
{\large A Conditional Reduction toward the Riemann Hypothesis via Spectral Microcluster Closure}\\[2pt]
{\large in the Connes--Consani--Moscovici Framework}}

\author{Lukas Geiger\\
\small Independent Researcher, Bernau, Germany\\
\small ORCID: \href{https://orcid.org/0009-0005-7296-1534}{0009-0005-7296-1534}}

\date{Preprint v1.6, 15 August 2026}

\begin{document}
\overfullrule=0pt
\maketitle

\begin{center}
\textit{Dedicated to Alain Connes, Caterina Consani, and Henri Moscovici, keepers of the spectral zoo}
\end{center}

\begin{abstract}
We present a spectral and numerical study of the
Connes--Consani--Moscovici (CCM) Fourier model for the Weil quadratic
form.  The Weil operator decomposes as a rank-one pole term plus an
arithmetic perturbation, and we prove an exact resolvent identity that
expresses the Riemann zeros as resonances where the direct pole
response matches the arithmetic resolvent response.

The main structural correction of the paper is the replacement of the
older fixed-tolerance cluster picture by a \emph{purely spectral
microcluster}: a low-energy eigenspace of the full Weil operator cut
out by an order-one outer gap.  This eliminates the circularity
objection to mass-based cluster definitions and leads to a clean
three-statement closure architecture for the CCM missing step MS2 in
the finite-dimensional Fourier model.  The three closure statements
are: scalar secular cancellation, projected Poisson quasimode control,
and a coercive complement outside the spectral microcluster.  At the
\emph{architectural} level each is established by an exact algebraic
identity; at the \emph{quantitative} level (uniform-in-$\lambda$
control of the constants $s_\lambda$, $p_\lambda$, $g_*$) each is
currently supported by stable numerics on the tested range
$\lambda \le 100$ but lacks a closed-form decay/lower-bound proof.
Combined via a standard quasimode-to-projector argument they yield
\[
  \|(I - P_{V_\lambda})\, k_\lambda\|
  \;\leq\; \frac{s_\lambda + p_\lambda}{g_*}
  \;\approx\; 6 \times 10^{-7}
  \quad\text{(on the tested range)}.
\]

Combined with two external CCM inputs --- the truncation-faithfulness
implication and the Hurwitz passage from finite approximants to the
limit, both stated in the source literature as proof
\emph{strategies} rather than completed theorems --- the present
argument yields a \emph{conditional reduction} of the Riemann
Hypothesis: RH follows if (i) the three closure statements hold
uniformly in $\lambda$ and (ii) the CCM external transfer is
rigorously established.  The earlier Abel--Stieltjes route is retained
as a diagnostic and structural precursor.  The result of this paper is
therefore best read as a contribution to the architecture and
numerical phenomenology of the CCM programme, not as a proof of RH.
\end{abstract}

\noindent\textbf{MSC 2020:} 11M26, 47A10, 47B35, 11M36\\
\textbf{Keywords:} Riemann Hypothesis, Weil explicit formula, noncommutative geometry,
spectral approximation, Abel summation, Cauchy interlacing

\newpage
\tableofcontents
\newpage

% =============================================================================
\section{Introduction: The Spectral Zoo}
\label{sec:intro}
% =============================================================================

Every attempt to prove the Riemann Hypothesis confronts the same tension:
the conjecture is a statement about the arithmetic of prime numbers,
yet every promising approach requires spectral or geometric machinery
far removed from elementary number theory.
Random matrix theory reveals the statistical DNA of the zeros
\cite{Montgomery1973,Odlyzko1987};
quantum chaos connects them to the semiclassical limit of a
hypothetical Hamiltonian \cite{BerryKeating1999};
the Selberg trace formula links them to geodesics on hyperbolic surfaces;
automorphic forms embed them in a representation-theoretic universe.
Each approach captures a different facet of the phenomenon.
Together, they form a \emph{spectral zoo}---a diverse collection of
mathematical creatures, each adapted to its own ecological niche,
each providing a partial view of the same underlying structure.

If this landscape is a zoo, then Alain Connes has served as its most
dedicated keeper.
For over three decades, Connes has systematically constructed the
noncommutative geometry framework that promises to unify these
disparate approaches into a single coherent habitat.
His programme began with the observation that the explicit formulae of
prime number theory---going back to Riemann himself and formalised by
Weil \cite{Weil1952}---have a natural spectral interpretation in the
language of operator algebras.
The zeros of the zeta function appear as an \emph{absorption spectrum}:
frequencies at which a certain operator absorbs incoming signals completely.

The decisive step came with the work of Connes, Consani, and Moscovici
\cite{CCM2025}, who provided a concrete Fourier basis
in which the Weil quadratic form becomes a finite-dimensional matrix
amenable to direct computation.
In this basis, the Riemann Hypothesis reduces to a single spectral
approximation property---the \emph{Main Spectral Hypothesis} (MS2)---stating
that a certain Poisson-type kernel $k_\lambda$ lives approximately
in the cluster eigenspace of the Weil operator $QW$.

This paper enters the zookeeper's enclosure.
We implement the CCM framework computationally, using arbitrary-precision
arithmetic (50~decimal digits) to build the Weil operator in dimensions
up to~$121$ ($N = 120$ Fourier modes), and we test its predicted
diagnostics numerically for the Riemann zeta function.
Our main contributions are:

\begin{enumerate}[label=(\roman*)]
\item An \textbf{exact resolvent identity} (Theorem~\ref{thm:resolvent})
  that expresses the Riemann zeros as resonances of a rank-one perturbation.
  (Analytically proved.)
\item The \textbf{$\eta$-parametrisation} (Theorem~\ref{thm:eta-equiv}),
  which reduces the Riemann Hypothesis to the spectral inequality
  $0 < \eta_j < 2$ for each regular mode.  The rank-one sign forcing
  behind $\eta_j>0$ is proved analytically; the finite microcluster
  correction and the upper bound $\eta_j<2$ are supported numerically and
  folded into the conditional endgame estimates below.
\item The \textbf{Abel--Stieltjes cancellation technique}
  (\cref{sec:abel}), an intermediate result that transforms the sign
  problem in the Cauchy sum into positive difference kernels and
  explains the earlier tail-vs-gap numerics.
\item The \textbf{Three-Lemma Microcluster Closure}
  (\cref{sec:endgame}): a spectral endgame for the CCM missing step
  based on scalar secular cancellation, projected Poisson quasimodes,
  and a coercive complement outside a purely spectral microcluster.
\end{enumerate}

The paper is structured as follows.
\Cref{sec:framework} introduces the CCM framework and its operator
decomposition.
\Cref{sec:mechanism} establishes the cancellation mechanism and the
resolvent identity.
\Cref{sec:proof-arch} develops the proof architecture via the
$\eta$-parametrisation.
\Cref{sec:abel} introduces the Abel--Stieltjes technique.
\Cref{sec:numerics} presents the numerical evidence.
\Cref{sec:endgame} develops the spectral-microcluster endgame that
reduces MS2 inside the finite-dimensional CCM model to three explicit
quantitative closure inputs.
\Cref{sec:RH} draws the consequence for the Riemann Hypothesis via
the standard CCM transfer and Hurwitz passage.
\Cref{sec:conclusion} concludes.

% =============================================================================
\section{The CCM Framework}
\label{sec:framework}
% =============================================================================

\subsection{The Weil explicit formula as a quadratic form}

The Weil explicit formula \cite{Weil1952} expresses a sum over the
non-trivial zeros $\rho$ of $\zeta(s)$ as a distributional identity
involving primes:
\begin{equation}
\label{eq:weil}
\sum_\rho \hat{h}(\rho - \tfrac{1}{2})
= \hat{h}(\tfrac{i}{2}) + \hat{h}(-\tfrac{i}{2})
- \sum_p \sum_{m=1}^\infty \frac{\log p}{p^{m/2}}
  \bigl[h(m\log p) + h(-m\log p)\bigr]
+ \text{(arch.)},
\end{equation}
where $h$ is a suitable test function and $\hat{h}$ its Fourier
transform.
Connes' key insight \cite{Connes1999} is that this formula defines
a \emph{quadratic form}
\begin{equation}
\label{eq:QW-def}
QW(h) = \text{pole terms} - \text{prime terms} + \text{archimedean terms},
\end{equation}
and the Riemann Hypothesis is equivalent to the statement that
$QW(h) \geq 0$ for all admissible test functions~$h$---the
\emph{Weil positivity criterion}.
This places the present route in the broader family of positivity
criteria for RH, alongside Li's sequence criterion \cite{Li1997}, while
using the Weil/CCM quadratic-form framework rather than Li coefficients.

\subsection{The Fourier basis of Connes--Consani--Moscovici}

Connes, Consani, and Moscovici \cite{CCM2025} introduce a specific
Fourier basis $\{e_1, \ldots, e_{2N+1}\}$ for the even part of the
Weil quadratic form.
In this basis, the operator $QW_{\mathrm{even}}$ becomes a
finite-dimensional Hermitian matrix of dimension $\dim = 2N+1$.
The bandwidth parameter $\lambda > 0$ controls the support of the
test function: larger $\lambda$ includes more primes and finer
spectral resolution.

The central structural result is the \emph{rank-one decomposition}:

\begin{proposition}[Rank-one decomposition]
\label{prop:rank1}
In the CCM Fourier basis, the Weil operator decomposes as
\begin{equation}
\label{eq:rank1}
QW_{\mathrm{even}} = \Ctilde \, |\utilde\rangle\langle\utilde| + H,
\end{equation}
where:
\begin{itemize}
\item $\Ctilde > 0$ is the $W_{02}$-eigenvalue (the contribution of
  the pole at $s = 1$),
\item $\utilde$ is the normalised eigenvector of the $W_{02}$-operator,
\item $H = -W_R - W'$ is the arithmetic operator, built from prime data.
\end{itemize}
\end{proposition}

The rank-one term carries the ``counting'' information (the pole of
$\zeta$ at $s = 1$, encoding the density of integers), while $H$ carries
the ``coding'' information (the primes, encoding multiplicative structure).

% added 2026-05-23: Foundation-Lock identity for the true minimal eigenvector
\begin{proposition}[Foundation-Lock: canonical form of the true minimal eigenvector]
\label{prop:foundation-lock}
Let $L := 2\log\lambda$ be the support length of the CCM interval
$[-L/2, L/2]$, and let $F(z) = \sum_{j} \xi_j / (z - 2\pi j / L)$ be the
rational function associated to the true minimal eigenvector $\xi = (\xi_j)$
of $QW_{\mathrm{even}}$ at bandwidth~$\lambda$ (CCM Eq.~5.25).  Define
\begin{equation}
\label{eq:xihat-def}
  \hat{\xi}(z) := \frac{2}{\sqrt{L}} \sin\!\bigl(\tfrac{zL}{2}\bigr)\, F(z).
\end{equation}
Then:
\begin{enumerate}[label=(\alph*)]
  \item $\hat{\xi}$ is an \emph{entire} function: the apparent poles of $F$ at
    $z = 2\pi k/L$ are cancelled by the zeros of $\sin(zL/2)$.
  \item The lattice evaluations satisfy
    $\hat{\xi}(2\pi k/L) = \sqrt{L}\,(-1)^k\,\xi_k$ for all $k$.
  \item $\hat{\xi}$ is even: $\hat{\xi}(-z) = \hat{\xi}(z)$.
\end{enumerate}
In particular, $\hat{\xi}$ and $F$ share the same non-trivial zeros
(with opposite parities: $F$ is odd, $\hat{\xi}$ is even).
\end{proposition}

\begin{proof}
\textit{Parts~(a) and~(b).}
Near $z_0 = 2\pi k / L$ the rational function $F$ has the Laurent expansion
\[
  F(z) = \frac{\xi_k}{z - z_0} + O(1),
\]
i.e.\ a simple pole with residue $\xi_k$.  For the sine factor, substituting
$z = z_0 + \varepsilon$ gives
\begin{align*}
  \sin\!\bigl(\tfrac{zL}{2}\bigr)
  &= \sin\!\bigl(\pi k + \tfrac{\varepsilon L}{2}\bigr)
   = (-1)^k \sin\!\bigl(\tfrac{\varepsilon L}{2}\bigr)
   = (-1)^k \tfrac{L}{2}\,\varepsilon + O(\varepsilon^3).
\end{align*}
Multiplying by the prefactor $(2/\sqrt{L})$:
\begin{align*}
  \hat{\xi}(z_0 + \varepsilon)
  &= \frac{2}{\sqrt{L}}\,(-1)^k\,\frac{L}{2}\,\varepsilon
     \cdot \frac{\xi_k}{\varepsilon} + O(\varepsilon)
   = \sqrt{L}\,(-1)^k\,\xi_k + O(\varepsilon).
\end{align*}
Since the $1/\varepsilon$ singularity cancels exactly, $\hat{\xi}$ extends
analytically to $z_0$---proving~(a).  Taking $\varepsilon \to 0$ yields
$\hat{\xi}(2\pi k/L) = \sqrt{L}\,(-1)^k\,\xi_k$---proving~(b).

\textit{Part~(c).}
$\sin((-z)L/2) = -\sin(zL/2)$ and $F(-z) = -F(z)$ (both are odd); their
product $\hat{\xi}(-z) = \hat{\xi}(z)$ is therefore even.
\end{proof}

\begin{remark}[Relation to the educated guess $k_\lambda$]
\label{rem:xihat-vs-klambda}
The Proposition concerns $\hat{\xi}$, the Fourier transform of the
\emph{true} minimal eigenvector $\xi$, not the educated guess $k_\lambda$
used throughout the remainder of this paper.  Whether $k_\lambda\approx\xi$
(i.e.\ whether MS2 holds) is the central question; the Foundation-Lock
identity \eqref{eq:xihat-def} gives $\hat{\xi}$ a canonical representation that is
independent of this approximation.  The identity has been verified
numerically at $\lambda=5, 7$ at \texttt{dps}$=50$, with maximum
deviation $\approx 10^{-52}$.
\end{remark}

\subsection{The Poisson kernel and MS2}

The CCM framework provides a distinguished vector, the
\emph{Poisson kernel}
\begin{equation}
\label{eq:poisson}
k_\lambda(u) = u^{1/2} \lambda^{1/2} \sum_{n=1}^\infty h(\lambda n u),
\end{equation}
which serves as an educated guess for the ground state of $QW$.
In the Fourier basis, $k_\lambda$ has coordinates $k_j = \ip{e_j}{k_\lambda}$.

\begin{definition}[Main Spectral Hypothesis]
\label{def:MS2}
We say \emph{MS2 holds at bandwidth $\lambda$} if the Poisson kernel
$k_\lambda$ lies approximately in the cluster eigenspace of $QW_{\mathrm{even}}$:
\begin{equation}
\label{eq:MS2}
\norm{P_\perp k_\lambda}^2 \to 0 \quad \text{as } \lambda \to \infty,
\end{equation}
where $P_\perp = I - P_{\Cl}$ projects onto the complement of the
cluster eigenspace (the eigenspace of the minimal eigenvalue $\wmin$).
\end{definition}

\begin{remark}
The Weil positivity $QW \geq 0$ (which is equivalent to RH) follows
from MS2 by a variational argument: if $k_\lambda$ lies in the cluster
eigenspace, then $QW(k_\lambda) = \wmin \norm{k_\lambda}^2 + O(\varepsilon_\lambda)$,
and the ratio $\wmin / \Ctilde \to 0$ ensures positivity in the limit.
\end{remark}

\subsection{Spectral data}

We denote by $A = QW_{\mathrm{even}}$ the full operator with
eigenvalues $w_a$ and eigenvectors $q_a$, and by $T = P_0 H P_0$
the projected arithmetic operator (where $P_0$ projects onto the
nullspace of $W_{02}$) with eigenvalues $t_j$ and eigenvectors $e_j$.
We write:
\begin{align}
\alpha &= \ip{\utilde}{k_\lambda}, &
\alpha_a &= \ip{\utilde}{q_a}, &
f_j &= \ip{e_j}{P_0 H \utilde}, \\
g_j &= \ip{e_j}{P_0 k_\lambda}, &
c_a &= \ip{q_a}{k_\lambda}, &
r_j &= \alpha f_j + (t_j - \wmin) g_j. \notag
\end{align}

The quantity $r_j$ is the \emph{residual}: it measures how well the
Poisson kernel is approximated by the cluster resolvent response at
mode~$j$.

% =============================================================================
\section{The Cancellation Mechanism}
\label{sec:mechanism}
% =============================================================================

\subsection{Fifteen-digit cancellation}

The CCM framework reveals a striking numerical phenomenon:
at each Riemann zero $\gamma_j$, the evaluation functional $F(\gamma)
= \ip{e(\gamma)}{q_k}$ (where $e(\gamma)$ is the evaluation vector
at the zero and $q_k$ a cluster eigenvector) exhibits cancellation
to~$15$ decimal places:

\begin{equation}
\label{eq:cancel}
F(\gamma_1) = \underbrace{+0.0387\ldots}_{\text{pole term}}
+ \underbrace{(-0.0387\ldots)}_{\text{nullspace term}}
= 1.45 \times 10^{-17}.
\end{equation}

At non-zeros, $F(z) \sim 0.3$--$2.6$---a separation factor of~$10^{13}$.
This cancellation is not accidental: it is enforced by the eigenvalue
equation $A q_k = w_k q_k$, which couples the pole and arithmetic
contributions in a rigid algebraic relation.

\subsection{The resolvent identity}

The algebraic mechanism behind the cancellation is captured by:

\begin{theorem}[Resolvent identity]
\label{thm:resolvent}
Let $q$ be a cluster eigenvector with $A q = \wmin q$ and $\beta = \ip{\utilde}{q} \neq 0$.
Then for any evaluation vector $e(\gamma)$:
\begin{equation}
\label{eq:resolvent}
F(\gamma) = \beta \cdot \bigl[\alpha(\gamma) - R(\gamma, \wmin)\bigr],
\end{equation}
where
\begin{align}
\alpha(\gamma) &= \ip{\utilde}{e(\gamma)}
  && \text{(direct pole response)}, \\
R(\gamma, w) &= \ip{P_0 e(\gamma)}{(T - w I)^{-1} P_0 H \utilde}
  && \text{(arithmetic resolvent response)}.
\end{align}
\end{theorem}

\begin{proof}
From $A q = \wmin q$ and $A = \Ctilde |\utilde\rangle\langle\utilde| + H$,
projecting onto $\ker W_{02}$ gives
$(T - \wmin I) P_0 q = -\beta P_0 H \utilde$.
Since $T - \wmin I$ is invertible on $\ker W_{02}$
(by the interlacing property $\wmin < t_j$ for all~$j$), we solve:
$P_0 q = -\beta (T - \wmin I)^{-1} P_0 H \utilde$.
Substituting into $F(\gamma) = \beta \alpha(\gamma) + \ip{P_0 e(\gamma)}{P_0 q}$
yields~\eqref{eq:resolvent}.
\end{proof}

\begin{corollary}[Zeros as resonances]
$F(\gamma) = 0$ if and only if $\alpha(\gamma) = R(\gamma, \wmin)$:
the Riemann zeros are \emph{resonances} at which the incoming pole
signal is perfectly absorbed by the arithmetic resolvent.
\end{corollary}

The resolvent identity has been verified numerically to~$10^{-13}$
relative accuracy across all cluster eigenvectors and the first five
Riemann zeros (\cref{tab:resolvent}).

% =============================================================================
\section{The Proof Architecture}
\label{sec:proof-arch}
% =============================================================================

\subsection{The \texorpdfstring{$\eta$}{eta}-parametrisation}

The key algebraic step is to decompose the Poisson kernel
$k_\lambda = \alpha \utilde + g$ with $g = P_0 k_\lambda$,
and to express the residual at each mode~$j$ as:
\begin{equation}
\label{eq:residual}
r_j = \alpha f_j + (t_j - \wmin) g_j.
\end{equation}

\begin{definition}[$\eta$-parameter]
For each regular mode $j \in J_{\reg}$ (modes with $|f_j| > 0$), define
\begin{equation}
\label{eq:eta-def}
\eta_j := -\frac{(t_j - \wmin)\, g_j}{\alpha\, f_j}.
\end{equation}
\end{definition}

The residual then takes the form $r_j = \alpha f_j (1 - \eta_j)$,
and the central equivalence is:

\begin{theorem}[$\eta$-equivalence]
\label{thm:eta-equiv}
The following are equivalent:
\begin{enumerate}[label=(\alph*)]
\item $|r_j| < \alpha |f_j|$ for all $j \in J_{\reg}$
  \hfill (modeweise comparison),
\item $\eta_j \in (0, 2)$ for all $j \in J_{\reg}$
  \hfill ($\eta$-bound),
\item $|g_j^\perp| < \frac{\alpha_{\Cl}\, |f_j|}{t_j - \wmin}$
  for all $j \in J_{\reg}$
  \hfill (off-cluster stability).
\end{enumerate}
Moreover, each of these implies $\varphi_j = g_j / f_j < 0$ for all
$j \in J_{\reg}$.
\end{theorem}

\begin{proof}
The equivalence (a) $\Leftrightarrow$ (b) follows directly from
$r_j = \alpha f_j(1 - \eta_j)$ and $|1 - \eta_j| < 1 \Leftrightarrow \eta_j \in (0,2)$.
For (b) $\Leftrightarrow$ (c), decompose $g = g^{\Cl} + g^\perp$ where
$g_j^{\Cl} = -\alpha_{\Cl} f_j / (t_j - \wmin)$
(from the cluster resolvent, Lemma~\ref{lem:cluster-resolvent}).
Then $\eta_j = \alpha_{\Cl}/\alpha - (t_j - \wmin) g_j^\perp / (\alpha f_j)$,
and the condition $\eta_j \in (0,2)$ reduces to (c) since
$\alpha_{\Cl}/\alpha \to 1$.
\end{proof}

\subsection{From \texorpdfstring{$\eta$}{eta}-positivity to MS2}

The remaining bridge from the modewise inequality to MS2 can be written
as a Stieltjes route inside the present paper.  The algebraic sign
statement is exact; the actual MS2 conclusion still requires the uniform
off-cluster decay estimates isolated later in the endgame section.

\begin{proposition}[From $\eta$-positivity to Stieltjes control]
\label{prop:eta-to-ms2}
Assume $\eta_j \in (0,2)$ for every $j \in J_{\reg}$, and write
\[
  \varphi_j := \frac{g_j}{f_j},
  \qquad
  \nu_{\lambda,j} := f_j g_j.
\]
Then $\varphi_j < 0$ and $\nu_{\lambda,j} \le 0$ on $J_{\reg}$.  Hence
the regular part of the transfer function
\[
  m_\lambda^{\reg}(w)
  :=
  \sum_{j \in J_{\reg}} \frac{\nu_{\lambda,j}}{t_j-w}
  =
  - \sum_{j \in J_{\reg}} \frac{\mu_{\lambda,j}}{t_j-w},
  \qquad
  \mu_{\lambda,j} := - \nu_{\lambda,j} \ge 0,
\]
is a negative discrete Stieltjes transform.  In particular:
\begin{enumerate}[label=(\roman*)]
\item $m_\lambda^{\reg}$ is monotone decreasing between regular poles;
\item for $w,w_\ast$ in the same regular pole-free interval,
\[
  \bigl|m_\lambda^{\reg}(w)-m_\lambda^{\reg}(w_\ast)\bigr|
  \le
  |w-w_\ast|
  \sum_{j\in J_{\reg}}
  \frac{\mu_{\lambda,j}}{|t_j-w|\,|t_j-w_\ast|};
\]
\item if the same control is available uniformly in the off-cluster
  pole-free regime, then the off-cluster coefficients
  $c_a = \beta_a\delta_a$ inherit the weighted flatness bound needed for
  MS2, so that
\[
  \sum_{a\notin \Cl} |c_a|^2 \to 0
  \quad\Longrightarrow\quad
  \norm{P_\perp k_\lambda}^2 \to 0.
\]
\end{enumerate}
Consequently, the $\eta$-bound identifies a sufficient Stieltjes route to
MS2; the remaining quantitative content is the uniform off-cluster decay
estimate, supplied only conditionally in \cref{sec:endgame}.
\end{proposition}

\begin{proof}
From $\eta_j \in (0,2)$ and
$\eta_j = -(t_j-\wmin)g_j/(\alpha f_j)$ with $t_j-\wmin>0$ and
$\alpha>0$, we obtain $\varphi_j = g_j/f_j < 0$.  Therefore
$\nu_{\lambda,j} = f_j^2 \varphi_j \le 0$, which yields the Stieltjes
representation for $m_\lambda^{\reg}$.  Monotonicity follows by
termwise differentiation:
\[
  \frac{d}{dw} m_\lambda^{\reg}(w)
  =
  -\sum_{j\in J_{\reg}}
  \frac{\mu_{\lambda,j}}{(t_j-w)^2}
  \le 0.
\]
The flatness estimate is the elementary identity
\[
  \frac{1}{t_j-w} - \frac{1}{t_j-w_\ast}
  =
  \frac{w-w_\ast}{(t_j-w)(t_j-w_\ast)}
\]
summed against the positive weights $\mu_{\lambda,j}$.  Under the
uniform separation and summability hypotheses stated later as the
endgame inputs, evaluating this regular Stieltjes control at
off-cluster eigenvalues produces the weighted flatness estimate for
$\delta_a := \alpha_\lambda-m_\lambda(w_a)$, and hence for the
coefficients $c_a = \beta_a\delta_a$.  This is the off-cluster decay
mechanism behind MS2; the uniform estimate itself is not proved in this
subsection. \qedhere
\end{proof}

\subsection{The lower bound: Cauchy interlacing}

\begin{theorem}[Rank-one sign forcing for $\eta$]
\label{thm:eta-positive}
At the rank-one cluster level, $\eta_j=1$ for all
$j\in J_{\reg}$.  In the finite microcluster model, $\eta_j>0$ follows
whenever the multi-rank correction
$\delta_j:=1-\eta_j$ satisfies $|\delta_j|<1$.  This correction is
numerically small on the tested range and is one of the quantities
controlled conditionally by the Abel--Stieltjes and endgame bounds.
\end{theorem}

\begin{proof}
We argue via the secular equation for rank-one perturbations, which
forces the sign of the components $g_j = \ip{q_j}{g}$ in the
$T = P_0 H P_0$ eigenbasis without invoking any small-correction
approximation.

\emph{Step 1 (Secular equation).}  Since
$A = \Ctilde \, |\utilde\rangle\langle\utilde| + H$ is a rank-one
perturbation of $H$ with positive coupling $\Ctilde > 0$, the
eigenvalues $w$ of $A$ that are not eigenvalues of $T$ satisfy the
classical secular equation
\begin{equation}
\label{eq:secular}
1 + \Ctilde \sum_a \frac{|\ip{q_a}{\utilde}|^2}{t_a - w}
\;=\; 0,
\end{equation}
where the sum runs over the eigenpairs $(t_a, q_a)$ of $T$.  This is
the standard Bunch--Nielsen--Sorensen identity for rank-one updates of
symmetric operators \cite[Sec.~8.5]{GolubVanLoan2013}.

\emph{Step 2 (Cauchy interlacing.)}  Equation~\eqref{eq:secular} forces
strict interlacing $\wmin < t_1 < w_1 < t_2 < w_2 < \cdots$, where $w_k$
are the non-cluster eigenvalues of $A$.  In particular $t_j - \wmin > 0$
for every $j \in J_{\reg}$.

\emph{Step 3 (Component formula and sign-forcing).}  The eigenvector
equation $A g = \wmin g$ projected onto $q_j$ gives
\begin{equation}
\label{eq:eigvec-projection}
t_j\, g_j \;+\; \Ctilde\, \alpha_j\, \ip{\utilde}{g} \;=\; \wmin\, g_j,
\quad\text{i.e.}\quad
g_j \;=\; -\frac{\Ctilde\, \alpha_j\, \ip{\utilde}{g}}{t_j - \wmin},
\end{equation}
where $\alpha_j = \ip{q_j}{\utilde}$.  Define
$\beta := \Ctilde\, \ip{\utilde}{g}$; in our cluster normalisation
$\beta = -\alpha_{\Cl}\,\ip{q_j}{f}/\alpha$ at the cluster level and is
of constant sign across $j \in J_{\reg}$ by the structure of the
cluster resolvent in \cref{lem:cluster-resolvent}.  Identifying
$f_j = \alpha\,\alpha_j / \beta$ (up to the overall sign of $\beta$),
\eqref{eq:eigvec-projection} reads
\[
g_j \;=\; -\frac{\beta\, \alpha_j}{t_j - \wmin}
\;=\; -\frac{\alpha\, f_j}{t_j - \wmin}.
\]
Since $t_j - \wmin > 0$ by Step~2, $\sgn(g_j) = -\sgn(f_j)$ exactly
(not approximately).  No off-cluster correction term enters this
identity; it is a consequence of the eigenvector equation and
interlacing alone.

\emph{Step 4 (Rank-one conclusion).}  By definition
$\eta_j = -(t_j - \wmin)\,g_j / (\alpha\, f_j)$, so
\[
\eta_j = -\frac{(t_j - \wmin)}{\alpha\, f_j}
        \cdot \biggl(-\frac{\alpha\, f_j}{t_j - \wmin}\biggr)
      = 1 > 0.
\]
This proves the rank-one sign-forcing claim.  In the finite microcluster
model the deviation of $\eta_j$ from $1$ measures precisely the cluster
\emph{multi-rank} correction (multiple cluster eigenvectors with
distinct $w_a \approx \wmin$ instead of an exact rank-one cluster).  It
is the $\delta_j = 1 - \eta_j$ analysed in \cref{sec:abel}; positivity
of the finite $\eta_j$ follows once $|\delta_j|<1$, while the uniform
control of this correction belongs to the conditional endgame estimates.
\end{proof}

\begin{remark}[On the previous proof sketch]
\label{rem:eta-positive-history}
Earlier drafts of this paper argued $\eta_j > 0$ via a near-cluster
approximation $g_j \approx g_j^{\Cl}$.  The secular-equation argument
above is preferable: it eliminates the approximation step and reduces
the sign-forcing to standard rank-one perturbation theory.  The result
$\eta_j = 1$ at the rank-one cluster level is consistent with the
numerical observation $|1 - \eta_j| < 10^{-2}$ uniformly on the tested
range; deviations of $\eta_j$ from $1$ reflect the multi-cluster
microstructure, not a sign failure.
\end{remark}

\subsection{The upper bound: the historical frontier}

Historically, the bound $\eta_j < 2$ was the decisive unresolved local
inequality.  Numerically, $\eta_j \approx 1$ with $|1-\eta_j|<10^{-2}$
for all tested modes (\cref{tab:eta}), so the margin to~$2$ is
macroscopic.  In the present paper, however, the final closure will no
longer proceed by proving $\eta_j<2$ directly.  Instead, \cref{sec:endgame}
packages the same mechanism into the spectral-microcluster three-lemma
argument, while the $\eta$-formalism remains the clearest local
coordinate system for the cancellation structure.

The remainder of this paper develops a technique that reduces
$\eta_j < 2$ to a concrete, analytically tractable inequality.

\begin{lemma}[Cluster resolvent]
\label{lem:cluster-resolvent}
For each cluster eigenvector $q_a$ with $A q_a = \wmin q_a$
and $\alpha_a = \ip{\utilde}{q_a}$, the cluster projection of~$g$
in the $T$-eigenbasis is:
\begin{equation}
g_j^{(a)} = -\frac{\alpha_a f_j}{t_j - \wmin},
\qquad
g_j^{\Cl} = \sum_{a \in \Cl} c_a g_j^{(a)}
= -\frac{\alpha_{\Cl} f_j}{t_j - \wmin},
\end{equation}
where $\alpha_{\Cl} = \sum_{a \in \Cl} c_a \alpha_a$.
\end{lemma}

% =============================================================================
\section{Abel--Stieltjes Cancellation}
\label{sec:abel}
% =============================================================================

\subsection{Why triangle inequalities fail}

The off-cluster defect admits the spectral decomposition
\begin{equation}
\label{eq:delta-cauchy}
\delta_j := 1 - \eta_j
= -\frac{t_j - \wmin}{\alpha}
\sum_{a \notin \Cl} \frac{c_a \alpha_a}{t_j - w_a}
+ \frac{\alpha - \alpha_{\Cl}}{\alpha}.
\end{equation}

The natural impulse is to apply the triangle inequality:
$|\delta_j| \leq (M_{\off}/\alpha) \cdot (t_j - \wmin) / \min_a |t_j - w_a|$,
where $M_{\off} = \sum_{a \notin \Cl} |c_a \alpha_a|$.
This bound \emph{diverges}: the mass ratio $M_{\off}/\alpha$ decreases as
$\lambda^{-1.8}$, but the geometric ratio $\max[(t_j-\wmin)/\min|t_j-w_a|]$
grows as $\lambda^{+7}$, producing a net divergence.

The failure is structural.
The Cauchy sum in~\eqref{eq:delta-cauchy} has a dominant pole at
each mode (the nearest $w_a$), which is almost exactly compensated
by the neighbouring pole on the opposite side (Cauchy interlacing).
The triangle inequality destroys this cancellation.
Numerically, the \emph{tightness}
$|\delta_j|/\text{bound} \approx 10^{-6}$---the bound is a million times
too large.

\subsection{The Abel transformation}

The correct approach is \emph{summation by parts} (Abel transformation),
which groups the Cauchy sum into differences that respect the sign
structure.

\begin{definition}[Mass coefficients]
Define $b_a := c_a \alpha_a / \alpha > 0$ for $a \notin \Cl$.
Order the off-cluster eigenvalues as $w_{(1)} < w_{(2)} < \cdots < w_{(n_{\off})}$
with corresponding masses $b_{(1)}, \ldots, b_{(n_{\off})}$.
Write:
\begin{equation}
B_m = \sum_{k=1}^{m} b_{(k)}, \qquad
C_{m+1} = \sum_{k=m+1}^{n_{\off}} b_{(k)} = B_{\mathrm{tot}} - B_m.
\end{equation}
\end{definition}

For a mode $j$ with $w_{(m)} < t_j < w_{(m+1)}$ (the \emph{interlacing
position}), Abel summation decomposes the Cauchy sum as:

\begin{proposition}[Abel decomposition]
\label{prop:abel}
\begin{equation}
\label{eq:abel}
\sum_{a \notin \Cl} \frac{b_a}{t_j - w_a}
= \underbrace{\frac{B_m}{t_j - w_{(m)}} + \frac{C_{m+1}}{w_{(m+1)} - t_j}}_{\text{principal pair (boundary)}}
- \underbrace{\sum_{\ell} K_\ell}_{\text{Abel kernels (difference)}},
\end{equation}
where the Abel kernels
\begin{equation}
K_\ell = B_{(\ell)} \cdot \frac{w_{(\ell+1)} - w_{(\ell)}}{(t_j - w_{(\ell)})(t_j - w_{(\ell+1)})} > 0
\end{equation}
are \textbf{positive} for all terms on the same side of~$t_j$.
\end{proposition}

The positivity of the Abel kernels is the crucial advance:
it means the boundary terms \emph{overestimate} the actual sum,
and the difference kernels provide a \emph{controlled correction}
with known sign.

\begin{remark}[Principal pair balance]
Numerically, the principal pair contributes exactly $\sim 50\%$
of the total Abel bound across all modes and all~$\lambda$.
This is not a coincidence: the Abel identity
``actual $=$ boundary $-$ kernels'' implies that when $|\text{actual}|
\ll \text{boundary}$ (high cancellation), the kernels nearly equal
the boundary---hence PP $\approx 50\%$.
\end{remark}

\subsection{The Tail-vs-Gap Lemma}

Since the Abel kernels are positive, the actual sum is bounded by
the boundary terms alone:
\begin{equation}
|\delta_j| \leq (t_j - \wmin) \cdot
\left(\frac{B_m}{t_j - w_{(m)}} + \frac{C_{m+1}}{w_{(m+1)} - t_j}\right).
\end{equation}

The left boundary term $B_m \cdot d / \text{gap}_L$ is controlled by
the fact that $B_m$ is bounded while gap$_L$ stays bounded away from zero
for bulk modes.
The right boundary term $C_{m+1} \cdot d / \text{gap}_R$ is the
critical one: it pits the \emph{mass tail} $C_{m+1}$ against the
\emph{interlacing gap} $w_{(m+1)} - t_j$.

\begin{conjecture}[Tail-vs-Gap conjecture]
\label{conj:tvg}
For all regular interlacing intervals $w_{(m)} < t_j < w_{(m+1)}$:
\begin{equation}
\label{eq:tvg}
\frac{(t_j - \wmin) \cdot C_{m+1}}{w_{(m+1)} - t_j} \leq \theta < 1,
\end{equation}
where $\theta$ is a universal constant independent of $\lambda$ and~$j$.
\end{conjecture}

The full implication chain is:

\begin{equation}
\label{eq:chain}
\boxed{
\text{Tail-vs-Gap}
\!\implies\! |\delta_j| < 1
\!\implies\! \eta_j \in (0,2)
\!\implies\! \varphi_j < 0
\!\implies\! \text{MS2}
\!\implies\! \text{RH}.
}
\end{equation}

\subsection{Mass concentration: the key mechanism}

The Tail-vs-Gap Lemma holds numerically because of a strong
\emph{mass concentration} property: the coefficients $b_a$ are
overwhelmingly concentrated near the cluster.

\begin{definition}[Mass concentration]
The \emph{$p$-concentration index} is the smallest $k$ such that
$B_k / B_{\mathrm{tot}} \geq p$.
\end{definition}

\begin{table}[h]
\centering
\caption{Mass concentration across bandwidth parameters.}
\label{tab:mass}
\begin{tabular}{cccccc}
\toprule
$\lambda$ & $n_{\off}$ & 50\% in first & 90\% in first & 99\% in first \\
\midrule
3 & 26 & 1/26 (3.8\%) & 2/26 (7.7\%) & 3/26 (11.5\%) \\
5 & 38 & 2/38 (5.3\%) & 4/38 (10.5\%) & 12/38 (31.6\%) \\
7 & 47 & 2/47 (4.3\%) & 5/47 (10.6\%) & 13/47 (27.7\%) \\
\bottomrule
\end{tabular}
\end{table}

Ninety percent of the off-cluster mass is concentrated in the
first $\sim 10\%$ of eigenvalues.
This means that for modes deep in the bulk ($m$ large),
the tail mass $C_{m+1}$ is tiny, easily overcoming the shrinking
gap.

\subsection{Worst-mode anatomy}

The worst mode at $\lambda = 7$ illustrates the mechanism:
mode $j = 60$ at interlacing position $m = 28$, where:
\begin{align*}
B_{28} &= 4.93 \times 10^{-6} && (99.87\% \text{ of total mass left of } t_j), \\
C_{29} &= 6.35 \times 10^{-9} && (0.13\% \text{ right}), \\
\text{gap}_R &= 4.53 \times 10^{-7}, \\
\theta &= C_{29} \cdot d / \text{gap}_R = 0.059 < 1. \quad \checkmark
\end{align*}

The tail mass $C_{29}$ has decayed by a factor of~$800$ relative
to $B_{\mathrm{tot}}$, easily compensating the small gap.

\begin{remark}[Superseded by Three-Lemma route]
The Abel--Stieltjes analysis above provides strong numerical evidence
and reveals the cancellation mechanism, but reduces the problem to the
monolithic Tail-vs-Gap conjecture (\cref{conj:tvg}).
\Cref{sec:endgame} develops a strictly more informative conditional
proof that identifies the \emph{three individual bounds} whose
analytical proof would close the argument.
\end{remark}

% =============================================================================
\section{Numerical Evidence}
\label{sec:numerics}
% =============================================================================

All computations use \texttt{mpmath} at $50$~decimal digits
(\texttt{DPS}$=50$) with the \texttt{build\_operators} routine
from the CCM Fourier basis implementation.
The three configurations are:

\begin{center}
\begin{tabular}{cccccc}
\toprule
$\lambda$ & $N$ & dim & \#primes & cluster size & build time \\
\midrule
3 & 30 & 31 & 4 & 5 & $\sim 210$\,s \\
5 & 55 & 56 & 9 & 18 & $\sim 400$\,s \\
7 & 77 & 86 & 15 & 33 & $\sim 740$\,s \\
\bottomrule
\end{tabular}
\end{center}

\subsection{Resolvent identity verification}

\begin{table}[h]
\centering
\caption{Resolvent identity~\eqref{eq:resolvent} verification at $\lambda = 3$.}
\label{tab:resolvent}
\begin{tabular}{lc}
\toprule
Test & Result \\
\midrule
Coupling equation & verified to $10^{-15}$ \\
Identity at zeros ($|G(\gamma_j)|$) & $\approx 10^{-13}$ \\
Identity at non-zeros ($|G(z)|$) & $\approx 10^{-1}$ to $10^{0}$ \\
Separation factor & $10^{11}$ \\
Secular equation & verified to $10^{-16}$ \\
\bottomrule
\end{tabular}
\end{table}

\subsection{MS2: leakage scaling}

\begin{table}[h]
\centering
\caption{Out-of-cluster leakage with $N \propto \lambda^2$ scaling.}
\label{tab:leakage}
\begin{tabular}{ccccccc}
\toprule
$\lambda$ & $N$ & cluster/dim & leakage $\varepsilon$ &
  $\norm{k^\perp}$ & $|F({\gamma_1})|$ & $|F_{\Cl}(\gamma_1)|$ \\
\midrule
3 & 30 & 5/31 (16\%) & $1.3 \times 10^{-8}$ & $1.1 \times 10^{-4}$ &
  $1.6 \times 10^{-7}$ & $4.2 \times 10^{-18}$ \\
5 & 55 & 18/56 (32\%) & $3.3 \times 10^{-10}$ & $1.8 \times 10^{-5}$ &
  $2.3 \times 10^{-8}$ & $5.5 \times 10^{-19}$ \\
7 & 85 & 33/86 (38\%) & $1.0 \times 10^{-10}$ & $1.0 \times 10^{-5}$ &
  $5.0 \times 10^{-8}$ & $4.1 \times 10^{-19}$ \\
\bottomrule
\end{tabular}
\end{table}

The leakage scales as $\varepsilon \sim \lambda^{-3.7}$ on this finite
test range, providing numerical support for the MS2 concentration
diagnostic.
The cluster projection gives $|F_{\Cl}(\gamma_1)| \approx 10^{-18}$---even
better than the full Poisson kernel.

\subsection{The \texorpdfstring{$\eta$}{eta}-distribution}

\begin{table}[h]
\centering
\caption{Distribution of $\eta_j$ across regular modes.}
\label{tab:eta}
\begin{tabular}{ccccccc}
\toprule
$\lambda$ & modes & $\eta_j$ range & median $\eta_j$ &
  margin to 0 & margin to 2 & $\max|\delta_j|$ \\
\midrule
3 & 27 & $[0.9993, 1.0000]$ & 0.99999 & 0.999 & 1.000 &
  $6.2 \times 10^{-4}$ \\
5 & 40 & $[0.9906, 1.0001]$ & 0.99999 & 0.991 & 1.000 &
  $9.4 \times 10^{-3}$ \\
7 & 50 & $[0.9975, 1.0003]$ & 0.99998 & 0.997 & 1.000 &
  $2.5 \times 10^{-4}$ \\
\bottomrule
\end{tabular}
\end{table}

All $\eta_j$ values cluster tightly around~$1$, with margins to
the critical boundaries of~$0$ and~$2$ remaining large.
The margin to~$2$ is always $\geq 0.99$---three orders of magnitude
beyond what is needed.

\subsection{Abel--Stieltjes bound}

\begin{table}[h]
\centering
\caption{Abel--Stieltjes bound and Tail-vs-Gap inequality across all
configurations.}
\label{tab:abel}
\small
\resizebox{\textwidth}{!}{%
\begin{tabular}{cccccccc}
\toprule
$\lambda$ & modes & bound $< 1$ & max PP bound & max $C \cdot d / \text{gap}$
  & max $|\delta|$ & worst $j$ & PP\% \\
\midrule
3 & 27 & \textbf{27/27} & 0.0094 & 0.0017 &
  $6.2 \times 10^{-4}$ & 21 & 50.1\% \\
5 & 40 & \textbf{40/40} & 0.267 & 0.267 &
  $1.7 \times 10^{-4}$ & 24 & 50.0\% \\
7 & 50 & \textbf{50/50} & 0.060 & 0.059 &
  $2.5 \times 10^{-4}$ & 60 & 50.0\% \\
\bottomrule
\end{tabular}%
}
\end{table}

All $117$ modes across three bandwidth parameters pass the
Tail-vs-Gap inequality.
The bound is non-monotonic in~$\lambda$ (spike at $\lambda = 5$,
decrease at $\lambda = 7$), consistent with the ``travelling wave''
interpretation: the worst-case mode migrates through the spectrum
as $\lambda$ increases, but the tail mass always decays faster
than the gap shrinks.

% =============================================================================
\section{The Three-Lemma Endgame and Spectral Microclusters}
\label{sec:endgame}
% =============================================================================

The earlier sections isolate the cancellation mechanism locally
($\eta$-coordinates, resolvent identity, Abel decomposition).  We now
repackage the closing step in the form in which it actually survives the
route audit: not as a tail-vs-gap conjecture and not as a mass-based
cluster definition, but as a spectral-microcluster theorem built from
three separate lemmas.

\subsection{Spectral microclusters}

\begin{definition}[Spectral microcluster]
\label{def:geff}
Fix a pre-registered order-one waterline $g_0 > 0$.  The resonant
microcluster is
\[
  V_\lambda(g_0) := \operatorname{span}\{q_j : u_j < u_0 + g_0\},
\]
where $u_0 \le u_1 \le \cdots$ are the eigenvalues of the full Weil
operator $A$ and $q_j$ the corresponding eigenvectors.  When $g_0$ is
fixed, we write $V_\lambda:=V_\lambda(g_0)$.
\end{definition}

\begin{remark}[Non-circularity]
This definition is purely spectral: $V_\lambda(g_0)$ is determined by the
spectrum of $A$ alone and does not refer to the Poisson kernel $k_\lambda$
or to its spectral mass distribution.  The statement that
$k_\lambda$ is almost contained in a uniformly chosen waterline subspace is
the conclusion of the endgame, not part of the definition.  The
nontrivial closure question is whether such a waterline can be chosen
uniformly in $\lambda$ while the low-waterline modes and the exterior gap
remain controlled.
\end{remark}

\subsection{Truncation stability}

All statements of this section are made at fixed $(\lambda,N)$ in the
CCM Fourier model.  The point of the spectral definition above is that it
is stable under Galerkin refinement: the number of eigenvalues inside the
microcluster may change with $N$, but the argument only uses the
existence of an order-one outer gap from $u_0$ to $V_\lambda^\perp$.
This avoids the pathology of the older fixed-tolerance cluster cuts,
which sliced through the interior of the genuine low-energy packet and
therefore produced erratic pseudo-gaps of size $10^{-6}$.

\subsection{The three endgame statements}

The three statements below are the closing inputs to
\cref{thm:MS2-closure}.  Each is supported by a structural identity
(rank-one secular splitting, Poisson transfer formula, Cauchy
interlacing against the spectrum of $H$) and by numerical
verification on the tested range $\lambda \le 100$.  At the
\emph{architectural} level (algebraic identities) they are proven; at
the \emph{quantitative} level (uniform-in-$\lambda$ control of the
constants $s_\lambda$, $p_\lambda$, $g_*$) they currently rest on
empirical observation rather than on closed-form decay rates with
explicit constants.  We therefore state them as \emph{conjectural
closure statements} and refer the analytical refinement to
\cref{rem:endgame-status}.

\begin{conjecture}[Scalar secular cancellation]
\label{lem:scalar-secular}
Let
\[
  \mu_\lambda := \ip{\utilde}{(A-u_0)k_\lambda}.
\]
Then
\[
  |\mu_\lambda| \le s_\lambda,
  \qquad
  \frac{\mu_\lambda}{\alpha_\lambda}
  =
  (\bar w-u_0)+\frac{\ip{f_\lambda}{k_\perp}}{\alpha_\lambda},
\]
where the two terms on the right are individually of order $10^{-2}$ and
cancel to order $10^{-7}$.
\end{conjecture}

\begin{proof}[Architectural identity and numerical support]
The identity is the exact rank-one secular splitting from the
Poisson--Rayleigh analysis: it is an algebraic consequence of the
eigenvalue equation $Ak = \bar w k$ projected onto $\utilde$.  For the
true ground state of $A$ the two terms cancel identically; for the
Poisson kernel $k_\lambda$ they cancel up to the Galerkin truncation
error.  Numerically the cancellation factor is between $150\times$
and $230\times$ and shows no growth in $\lambda$ once $N/L$ is kept
fixed (\cref{tab:endgame}).  A closed-form decay rate $s_\lambda \to 0$
with explicit constants is conjectural; the natural candidate is
$s_\lambda = O(e^{-cN/L})$ from Galerkin truncation theory, but
\emph{uniform} control across $\lambda$ has not been derived in closed
form.
\end{proof}

\begin{conjecture}[Projected Poisson quasimode]
\label{lem:projected-poisson}
Let
\[
  h_\lambda := P_0(A-u_0)k_\lambda.
\]
Then
\[
  \norm{h_\lambda} \le p_\lambda,
  \qquad
  h_\lambda
  =
  \frac{1}{n_{L,N}}\,P_0\Pi_{L,N}\!\bigl(E_L^{\mathrm{bulk}} + B_L\bigr).
\]
\end{conjecture}

\begin{proof}[Architectural identity and numerical support]
The exact Poisson transfer formula is
\[
  h_\lambda
  = (n_{L,N})^{-1}\,P_0\Pi_{L,N}
    \bigl(E_L^{\mathrm{bulk}} + B_L\bigr).
\]
It follows from the bulk--boundary decomposition
$H_L K_L = H_\infty K + B_L$ together with
$k_\lambda = K_L / n_{L,N}$.  The bulk defect $E_L^{\mathrm{bulk}}$
is almost entirely aligned with $\utilde$, so its $P_0$-projection is
small; the boundary term $B_L$ is controlled by Poisson kernel decay
at the truncation boundary.  Numerically
$\norm{h_\lambda} \approx 2$--$3\times 10^{-6}$ uniformly on the tested
range (\cref{tab:endgame}).  A closed-form $p_\lambda \to 0$ with
explicit constants requires uniform quantitative control of both the
$P_0$-leakage and the boundary-decay rate; that control is
conjectural at present.
\end{proof}

\begin{conjecture}[Coercive complement, uniform-$g_*$]
\label{lem:coercive-complement}
There exists an order-one constant $g_* > 0$, independent of
$\lambda$, such that
\[
  \ip{(A-u_0)v}{v} \ge g_* \norm{v}^2
  \qquad\text{for all } v \in V_\lambda^\perp.
\]
Equivalently,
\[
  \inf \sigma(A|_{V_\lambda^\perp}) - u_0 \ge g_*.
\]
\end{conjecture}

\begin{proof}[Architectural identity and numerical support]
Once $V_\lambda$ is defined spectrally, the coercive estimate reduces
to a strictly positive uniform lower bound on $u_{J+1} - u_0$, where
$u_{J+1}$ is the first eigenvalue of $A$ outside $V_\lambda$.  The
rank-one interlacing picture relates this to the lower spectral gap of
the arithmetic operator $H$: the erratic
$10^{-6}$ spacings belong to the interior micro-splitting of the low
cluster, whereas the first genuine outer jump is separated by a
macroscopic gap that tracks the lower spectral gap of the arithmetic
operator $H$.
\end{proof}

\subsection{The closure theorem (conditional)}

\begin{theorem}[Conditional three-statement closure of MS2]
\label{thm:MS2-closure}
With $V_\lambda$ as in \cref{def:geff}, assume the three conjectural
closure statements \cref{lem:scalar-secular,lem:projected-poisson,lem:coercive-complement}.
Then
\begin{equation}
\label{eq:mass-bound}
  \norm{(I - P_{V_\lambda}) k_\lambda}
  \;\leq\; \frac{s_\lambda + p_\lambda}{g_*}.
\end{equation}
In particular, if the truncation passage is chosen so that
$s_\lambda + p_\lambda \to 0$ and $g_* \ge g_0 > 0$ is uniform in
$\lambda$, then MS2 follows.
\end{theorem}

\begin{remark}[Status of the three closure statements]
\label{rem:endgame-status}
\Cref{thm:MS2-closure} is conditional on three quantitative inputs whose
algebraic identities are established but whose uniform-in-$\lambda$
control is currently empirical:
\begin{enumerate}[label=(\roman*),leftmargin=*]
  \item \emph{Scalar secular cancellation
  \textup{(\cref{lem:scalar-secular})}.}  The identity
  $\mu_\lambda/\alpha = (\bar w - u_0) + \ip{f_\lambda}{k_\perp}/\alpha$
  is an algebraic consequence of the eigenvalue equation projected onto
  $\utilde$.  The bound $|\mu_\lambda| \le s_\lambda$ holds with
  $s_\lambda \approx 4$--$5 \times 10^{-7}$ on the tested range; a
  uniform-in-$\lambda$ decay rate $s_\lambda \to 0$ is conjectured but
  not derived in closed form.
  \item \emph{Projected Poisson quasimode.}
  See \cref{lem:projected-poisson}.  The factorisation
  $h_\lambda = (n_{L,N})^{-1} P_0 \Pi_{L,N}(E_L^{\mathrm{bulk}} + B_L)$
  is an algebraic identity from the bulk--boundary decomposition.  The
  bound $\norm{h_\lambda} \le p_\lambda$ holds numerically with
  $p_\lambda \approx 2$--$3 \times 10^{-6}$; uniform analytic control
  of both the $P_0$-leakage and the boundary-decay rate is conjectural.
  \item \emph{Coercive complement.}
  See \cref{lem:coercive-complement}.  This is the most
  delicate of the three.  Existence of a uniform order-one outer gap
  $g_* \ge g_0 > 0$ requires a quantitative lower bound on
  $u_{J+1} - u_0$ that does not degrade as $\lambda \to \infty$.  The
  rank-one interlacing picture reduces this to the lower spectral gap
  of the arithmetic operator $H$, but a uniform-in-$\lambda$ bound on
  the latter is not established here.  Numerically
  $g_* \in [5.05, 7.13]$ on the tested range
  $\lambda \in \{30, 50, 70, 100\}$ (\cref{tab:endgame}); whether this
  remains bounded below by an explicit positive constant for all
  $\lambda$ is the most consequential internal open question of the
  present approach.
\end{enumerate}
The closure \cref{thm:MS2-closure} is therefore a \emph{conditional
reduction} of MS2 to (i)--(iii), not an unconditional proof of MS2.
\end{remark}

\begin{proof}
Set
\[
  R_\lambda := (A-u_0)k_\lambda = \mu_\lambda \utilde + h_\lambda.
\]
By the two residual lemmas,
\[
  \norm{R_\lambda}
  \le |\mu_\lambda| + \norm{h_\lambda}
  \le s_\lambda + p_\lambda.
\]
Now decompose $k_\lambda = k_{\mathrm{cl}} + k_\perp$ with
$k_{\mathrm{cl}} \in V_\lambda$ and $k_\perp \in V_\lambda^\perp$.
Since $V_\lambda$ is an $A$-invariant spectral subspace,
\[
  \ip{k_\perp}{(A-u_0)k_{\mathrm{cl}}} = 0,
\]
and therefore
\[
  g_* \norm{k_\perp}^2
  \le
  \ip{k_\perp}{(A-u_0)k_\perp}
  =
  \ip{k_\perp}{R_\lambda}
  \le
  \norm{k_\perp}\,\norm{R_\lambda}.
\]
Dividing by $\norm{k_\perp}$ gives
\[
  \norm{k_\perp}
  \le
  \frac{\norm{R_\lambda}}{g_*}
  \le
  \frac{s_\lambda+p_\lambda}{g_*},
\]
which is exactly \eqref{eq:mass-bound}.
\end{proof}

\subsection{Numerical calibration}

\begin{table}[h]
\centering
\caption{Spectral-microcluster closure: numerical calibration of the
three endgame constants.}
\label{tab:endgame}
\begin{tabular}{cccccc}
\toprule
$\lambda$ & $|\mu_\lambda|$ & $\norm{h_\lambda}$ & $g_*$ &
  $(|\mu_\lambda|+\norm{h_\lambda})/g_*$ &
  $\norm{(I-P_{V_\lambda})k_\lambda}$ \\
\midrule
30 & $2.7 \times 10^{-7}$ & $2.67 \times 10^{-6}$ & 5.05 & $5.83 \times 10^{-7}$ &
  $4.41 \times 10^{-7}$ \\
40 & $3.0 \times 10^{-7}$ & $2.96 \times 10^{-6}$ & 5.89 & $5.53 \times 10^{-7}$ &
  $4.16 \times 10^{-7}$ \\
50 & $2.8 \times 10^{-7}$ & $2.86 \times 10^{-6}$ & 5.90 & $5.32 \times 10^{-7}$ &
  $4.09 \times 10^{-7}$ \\
75 & $3.1 \times 10^{-7}$ & $3.03 \times 10^{-6}$ & 7.13 & $4.68 \times 10^{-7}$ &
  $3.93 \times 10^{-7}$ \\
100 & $2.9 \times 10^{-7}$ & $2.95 \times 10^{-6}$ & 6.82 & $4.75 \times 10^{-7}$ &
  $3.98 \times 10^{-7}$ \\
\bottomrule
\end{tabular}
\end{table}

The closure bound is 75--84\% tight, which indicates that the
quasimode-to-projector step is close to optimal.  More importantly, the
table exposes the decisive structural fact: the true outer gap is not a
microscopic spacing but an order-one constant, while the residual is
already of size $10^{-6}$.

% added 2026-05-23: §E.18 converged empirics from EvenDom companion
\subsection{Supporting evidence from the EvenDom companion: converged criterion-object norms}
\label{subsec:evendom-empirics}

The Even-Dominance companion paper \cite{Geiger2026evendom} studies the
true minimal eigenvector $\xi_\lambda$ and its canonical form $\hat{\xi}$
(see \cref{prop:foundation-lock}) via high-precision diagonalisation.
A subtle convergence issue arises: owing to the near-degeneracy of the
even/odd eigenvalue split (even/odd gap $\sim 10^{-65}$--$10^{-84}$,
super-exponentially small), accurate eigenvector computation requires
$\mathtt{dps}_{\min}(\lambda) \approx 30$--$50 \cdot \lambda$ decimal
places.  The main Zookeeper computations at $\mathtt{dps}=50$ remain
valid for $\lambda \le 100$ (since those computations target $k_\lambda$,
not $\xi_\lambda$), but any direct computation of $\hat{\xi}_\lambda$ at
larger $\lambda$ requires higher precision.

The following table reports converged values (from \cite{Geiger2026evendom}~§E.18)
for the sup-norm and $L^2$-norm proxies that enter the (ii-a) condition:
\begin{equation}
  A_\lambda(\sigma) := \sup_{t\in\mathbb{R}} |\hat{\xi}_\lambda(\sigma+it)|
  / |\hat{\xi}_\lambda(0)|,
  \qquad
  B_\lambda(\sigma) := \|\hat{\xi}_\lambda(\sigma+i\,\cdot\,)\|_{L^2(\mathbb{R})}
  / |\hat{\xi}_\lambda(0)|.
\end{equation}

\begin{table}[h]
\centering
\caption{Converged criterion-object norms at $\sigma=0.4$
  (from \cite{Geiger2026evendom}~§E.18; dps\,$=250$--$420$).
  $A_\lambda(0.4)\approx 1.004$ flat across $\lambda$: the sup-form of
  (ii-a) is numerically consistent with a uniform strip bound, but not
  a proof of it.  $B_\lambda(0.4)/\lambda^{0.4}
  \approx 0.68$ plateau: $L^2$-growth is sub-generic.}
\label{tab:evendom-norms}
\begin{tabular}{ccccc}
\toprule
$\lambda$ & dps & $A_\lambda(0.4)$ & $B_\lambda(0.4)$ &
  $B_\lambda(0.4)/\lambda^{0.4}$ \\
\midrule
3  & 250 & 1.0033 & 1.151 & 0.74 \\
5  & 300 & 1.0037 & 1.340 & 0.70 \\
7  & 350 & 1.0038 & 1.505 & 0.69 \\
9  & 420 & 1.0039 & 1.649 & 0.68 \\
11 & 440 & 1.0039 & 1.778 & 0.68 \\
13 & 540 & 1.0039 & 1.895 & 0.68 \\
15 & 640 & 1.0039 & 2.003 & 0.68 \\
\bottomrule
\end{tabular}
\end{table}

The flatness of $A_\lambda(0.4)$ across $\lambda = 3, 5, 7, 9, 11, 13, 15$ is
numerical evidence in favour of the a-priori strip bound (ii-a):
the sup-form remains $O(1)$ with no visible $\lambda$-growth.  The
$B_\lambda/\lambda^{0.4}$ plateau near $0.68$ indicates that the $L^2$
mass is sub-generic (generic growth would be $\sim\lambda^{1/2}$).
These results complement and are consistent with the internal Zookeeper
numerics: the Poisson-kernel results at $\mathtt{dps}=50$ and the true
eigenvector results at $\mathtt{dps}\ge 250$ are sampling different objects
($k_\lambda$ vs.\ $\xi_\lambda$) but both point toward the same
a-priori-boundedness structure.

\paragraph{26 May 2026 guardrail.}
The EvenDom follow-up audits sharpen what can and cannot be imported into
the Zookeeper route. The strip-curvature refit now uses six points through
$\lambda=20$, with $\kappa(20)=1.527\times10^{-3}$ and conservative
saturation range $\kappa_\infty\in[1.6,1.8]\times10^{-3}$. This supports
the conditional numerical picture behind (ii-a), but the uniform
log-curvature estimate remains open. Separately, the GF1 Chebyshev-ladder
probe has a low-bandwidth exception and an unresolved high-$\lambda$
discriminator: $\lambda=5,7$ support the benchmark
$\mathrm{sLI}\sim(64\pi^2/3)\lambda/(NL)$, $\lambda=3$ behaves as a
separate branch with older $C_m$ scale near~$7$, and the single
$\lambda=11,N=200$ point gives $C_{\mathrm{eff}}\approx168$ until the
larger-$N$ sweep resolves it. For the Zookeeper programme this means that
GF1 should be treated only as a diagnostic analogy for
Low-Waterline/Slush residual control; the structural route remains the
C2ca.1--C2ca.5 microcluster closure with uniform constants,
not a direct transfer of the EvenDom ladder.

\subsection{Relationship to earlier routes}

The Abel--Stieltjes analysis of \cref{sec:abel} remains valuable as a
mechanism finder: it reveals why triangle inequalities fail and why the
off-cluster cancellation is highly structured.  But the final closure no
longer depends on the tail-vs-gap conjecture as its single bottleneck.
The conditional Three-Lemma route is
strictly sharper: it replaces one monolithic conjecture by a spectral
microcluster statement and two residual estimates, each directly visible
in the operator decomposition.

\subsection{A 2026 cross-route sharpening: the determinant channel}
\label{sec:det-channel-sharpening}

A 2026 cross-route diagnostic phase (recorded in the Landscape mother, Part~IV)
sharpens the spectral target of this route in three ways, none of which changes
the conditional-reduction status.

\begin{itemize}
  \item \textbf{Archimedean density, independently confirmed.} A non-circular
    discretization of the genuine Connes--Moscovici prolate operator
    (arXiv:2112.05500) reproduces Riemann's zero-counting function $N(E)$ to
    $\sim 3\%$ in the UV --- the archimedean half (leading
    $\tfrac12\log(y/2\pi)$). This is an independent, non-circular corroboration
    of the CCM density claim underlying the microcluster construction.
  \item \textbf{The target is the determinant channel, not bare eigenvalue
    convergence.} A residue calibration shows that eigenvalue convergence
    $\mathrm{Spec}(A_F)\to\{\gamma_n\}$ is \emph{not} sufficient: the object that
    must converge is the regularized determinant
    $\det_{\mathrm{reg}}(A_F-z)\to C\cdot\Xi(z)$, equivalently the Weyl/trace
    channel $M_F=-D_F'/D_F$ with residues equal to multiplicities (a cyclic
    channel with the wrong normalization misses the Euler anchor by $\sim 50\%$;
    an off-axis Poisson tomography catches it $215$--$405\times$). \emph{Consequence
    for this route:} truncation faithfulness and the MS2 transfer must control the
    spectral \emph{measure} (poles \emph{and} residues), not only the spectrum.
  \item \textbf{The relative-determinant bypass is refuted.} An internally
    proposed shortcut --- removing the upper-half-plane Euler pole comb by a
    relative determinant $-\partial_z\log[\det(A-z)/\det(A^0-z)]$ against a prolate
    background --- does not work: the prolate background is itself Herglotz
    (comb-free), so there is nothing to cancel, and the only comb-free alternative
    requires a self-adjoint prime factor, which is circular. The determinant
    channel as a \emph{target} (previous item) remains valid; only the
    \emph{relative} determinant as a \emph{bypass} of the continuation wall is
    refuted.
\end{itemize}

\begin{table}[h]
\centering
\small
\setlength{\tabcolsep}{3pt}
\begin{tabular}{p{0.18\textwidth}p{0.28\textwidth}p{0.25\textwidth}p{0.19\textwidth}}
\toprule
Channel item & Data or diagnostic anchor & Present status & Overclaim blocked \\
\midrule
Archimedean density & Non-circular prolate discretization reproduces $N(E)$ to $\sim 3\%$ in the UV. &
Supports the smooth density side of the CCM model. & Not evidence for the arithmetic zero positions. \\
Determinant target & Residue calibration of $M_F=-D_F'/D_F$; residues must match multiplicities. &
Required target for truncation faithfulness. & Bare eigenvalue convergence is not enough. \\
Relative-determinant bypass & Prolate background is comb-free; a self-adjoint prime factor would be circular. &
Refuted as a shortcut. & Does not remove the analytic-continuation wall. \\
Cartwright bookkeeping & Canonical pairing is naked $P_N$ against $C\cdot\Xi$; the comb is internal to $P_N$. &
Presentation target for the next Zookeeper companion version. & No MS2 or RH upgrade; only normalization hygiene. \\
\bottomrule
\end{tabular}
\caption{Determinant-channel status ledger. The table separates diagnostic support, open transfer targets,
and refuted shortcuts so that the determinant channel is not read as a new closure of MS2.}
\label{tab:det-channel-ledger}
\end{table}

\noindent
\Cref{tab:det-channel-ledger} summarizes the status separation. The missing
datum remains the analytic continuation across $\Re s=1$ (the global limit, where
the local data must assemble into $\zeta$): the prolate operator's poles sit
$O(1)$ off the zeros, so the \emph{positions} (the arithmetic fine-structure) are
not supplied by the smooth density alone. This sharpens the spectral target; it
is not a closure. No MS2 or RH claim.


% =============================================================================
\section{Conditional Reduction to the Riemann Hypothesis}
\label{sec:RH}
% =============================================================================

To turn the microcluster closure into the Riemann Hypothesis, two
external inputs from the Connes--Consani--Moscovici programme are
required.  Both are explicitly formulated by their authors as proof
\emph{strategies}, not as completed theorems: \cite{CCM2025} describes
itself as ``strategy toward a proof of the Riemann Hypothesis'' and
\cite{Connes2026} as ``outlining a potential proof strategy based on
establishing convergence of zeros from finite to infinite Euler
products''.  We therefore record them here as named conjectural
inputs and present \cref{thm:RH-conditional} below as a
\emph{conditional reduction}, not as a proof of RH.

\begin{conjecture}[CCM truncation faithfulness, after \cite{CCM2025,Connes2026}]
\label{conj:ccm-blackbox}
Assume MS2 holds along an admissible truncation regime.  Then the
corresponding even CCM approximants satisfy the real-zero property for
the truncated $\xi_\lambda$.
\end{conjecture}

\begin{hypothesis}[External Hurwitz/prolate-wave input, after Fact 6.4 of \cite{Connes2026}]
\label{conj:hurwitz-blackbox}
If the CCM approximants $\xi_\lambda$ have the real-zero property along
a sequence $\lambda_n \to \infty$, then the Fourier transforms
$\widehat{\xi_{\lambda_n}}$ converge to the Riemann $\Xi$-function
uniformly on closed substrips $|\Im(z)| \le \delta < 1/2$ at rate
$O(\lambda^{-1/2-\delta})$, and by Hurwitz' theorem on zeros of
locally uniform limits of entire functions of order one, the
real-zero property passes to $\Xi$.  Equivalently, the Riemann
Hypothesis follows from real-zero approximants along a divergent
sequence.
\end{hypothesis}

\paragraph{Status of this statement.}  As stated by Connes in
\cite[Fact~6.4 of §6.5]{Connes2026}, the Fourier convergence with
the stated rate is presented as an established consequence of the
classical Slepian--Pollak prolate spheroidal wave function theory
\cite[Theorem~1 of reference~47]{Connes2026}, not as a conjecture.
The ``strategy'' framing in the abstract of \cite{Connes2026} refers
to the overall RH-architecture (the combination of Fact~6.4 with
even dominance and the approximation quality $k_\lambda\approx\xi_\lambda$),
not to Fact~6.4 itself.  The remaining open analytical questions
identified by Connes in \cite[§6.6]{Connes2026} are: (i)~simple
even ground state of $QW_\lambda$ (even dominance), addressed in the
companion paper \cite{Geiger2026evendom}; (ii)~approximation
quality of the educated guess $k_\lambda$ versus the true ground
state $\xi_\lambda$, which is precisely the spectral microcluster
question MS2 addressed in \cref{sec:endgame} of the present paper.

% added 2026-05-23: §E.14 reduction of limit identification (ii-c) into the a-priori bound (ii-a)
\subsection{Reduction of limit identification to the a-priori strip bound}
\label{subsec:iic-to-iia}

The Hurwitz passage (\cref{conj:hurwitz-blackbox}) requires condition~(ii)
in Connes' programme: $\hat{\xi}_\lambda \to \Xi$ locally uniformly on a
strip.  Following the Normal-Family framework of \cite{Geiger2026evendom},
condition~(ii) decomposes into two parts:
\begin{itemize}
  \item[(ii-a)] The family $\{\hat{\xi}_\lambda\}$ is locally bounded
    in the strip $|\Im(z)| < 1/2$ (a priori bound, needed for Montel
    compactness).
  \item[(ii-c)] Every subsequential limit of $\{\hat{\xi}_\lambda\}$ in
    the strip equals $\Xi$ (limit identification).
\end{itemize}
A standard Montel--Vitali argument gives (ii-a)$+$(ii-c)$\Rightarrow$(ii).

\begin{proposition}[Limit identification under zero-multiset coincidence and normalization]
\label{prop:iic-absorbed}
Assume (ii-a), and assume \textup{(Z+)}: every subsequential locally
uniform limit $G$ has the same full zero multiset as $\Xi$,
with the same normalization $G(0)=\Xi(0)$.  Then that subsequential
limit is $\Xi$.  Thus, once the zero multiset and normalization are
already controlled, the only remaining freedom is the exponential-factor
growth control supplied by the a-priori bound.
\end{proposition}

\begin{proof}[Proof sketch (five steps, from \cite{Geiger2026evendom} §E.14)]
\emph{Step~1.}  Because $\hat{\xi}_\lambda$ is even by
\cref{prop:foundation-lock}(c), any subsequential locally uniform limit
$G$ inherits parity $G(-z)=G(z)$---without using~(Z+).

\emph{Step~2.}  Under~(Z+) and local uniform convergence, the full zero
multiset of $G$ equals that of $\Xi$.  Since both functions are entire of
order~$1$, the Hadamard product gives
$G(z) = \Xi(z)\,e^{h(z)}$ for some entire $h$.

\emph{Step~3.}  Parity of both $G$ and $\Xi$ forces $h(-z)=h(z)$, i.e.\
$h$ is even; the real-axis restriction forces $h$ real-valued on~$\mathbb{R}$.

\emph{Step~4.}  The order-one growth bound inherited from the CCM setup
(with (ii-a) on the strip) forces $h$ to be a polynomial of degree
$\le 1$; combined with evenness, $h$ is a constant $h(z) = c_0$.

\emph{Step~5.}  The normalization $G(0)=\Xi(0)$ fixes $c_0 = 0$.
Hence $G = \Xi$.
\end{proof}

\begin{remark}[Status]
Assumption~(Z+) is itself a strong part of the limiting problem.  The
Proposition therefore does not eliminate Connes' condition~(ii); it only
records that, after zero-multiset coincidence and normalization have
already been supplied by another argument, the residual ambiguity is the
standard exponential factor, which is killed by order and normalization
control.  The open analytic target remains a rigorous a-priori strip
bound from $QW_N$-coercivity together with a non-circular source of
zero-multiset control.
\end{remark}

\begin{theorem}[Conditional reduction of RH via microcluster closure]
\label{thm:RH-conditional}
Assume \cref{conj:ccm-blackbox}, \cref{conj:hurwitz-blackbox}, and the
three endgame statements of \cref{thm:MS2-closure}.  Then the Riemann
Hypothesis holds.
\end{theorem}

\begin{proof}
\Cref{thm:MS2-closure} gives MS2 from the three endgame statements
(scalar secular cancellation, projected Poisson control, coercive
complement).  By \cref{conj:ccm-blackbox}, MS2 forces the real-zero
property for the CCM approximants.  By \cref{conj:hurwitz-blackbox},
this real-zero property passes to the completed zeta function in the
limit.  Therefore the Riemann Hypothesis holds.
\end{proof}

\paragraph{Honest statement of conditioning.}  The above is
\emph{not} a proof of the Riemann Hypothesis.  After the
clarifications above, the conditioning consists of:
\begin{itemize}[leftmargin=*]
  \item \emph{External imported input:} the Slepian--Pollak prolate
    wave convergence statement used in \cite[Fact 6.4 / Theorem~1 of
    Ref.~47]{Connes2026}, recorded here as
    \cref{conj:hurwitz-blackbox} and not reproved in this paper; and
    the Connes--van Suijlekom quadratic-form real-zero criterion
    \cite{ConnesvS2025}.
  \item \emph{Internal open inputs:} the three endgame statements of
    \cref{sec:endgame} (scalar secular cancellation, projected Poisson
    quasimode, coercive complement), which together close the spectral
    microcluster mechanism for MS2.  These remain as conjectural
    closure statements supported by an architectural identity plus
    numerical stability on the tested range.
  \item \emph{Internal open input (Connes' own remaining step):}
    simple even ground state of $QW_\lambda$ (Connes 2026 §6.6 item
    (i)), addressed in the companion paper \cite{Geiger2026evendom}.
\end{itemize}
The present paper addresses the MS2 question; the Even Dominance
question is handled in \cite{Geiger2026evendom}.  Both papers
together attempt to discharge exactly the two ``remaining steps'' that
Connes identifies in \cite[§6.6]{Connes2026}.

\begin{remark}[Independent route]
A separate conditional reduction of the Riemann Hypothesis via
\emph{even dominance} of the Weil quadratic form---using spectral
positivity methods rather than the microcluster mechanism---is
presented in \cite{Geiger2026evendom}, where it is itself conditional
on a quantitative variational gap conjecture for the odd-sector minimum
eigenvalue plus the same Connes program.  The two routes share the
CCM operator framework but diverge in the closing argument: even
dominance exploits the sign structure of the full quadratic form,
while the present paper isolates the quasimode concentration inside a
spectral microcluster.
\end{remark}

% added 2026-05-23: err-collapse criterion from EvenDom companion
\begin{remark}[err-collapse criterion from the companion paper]
\label{rem:err-collapse}
The companion paper \cite{Geiger2026evendom} formulates a second
reduction criterion, proved there under its stated metric and
normalization assumptions.  Define
$\mathrm{err}_\lambda := \sup_j |z_j^{(\lambda)} - \gamma_j|$ where
$z_j^{(\lambda)}$ are the non-trivial zeros of $F_{\xi_\lambda}$ (the
rational function of the true minimal eigenvector, CCM~Eq.~5.25) and
$\gamma_j$ the Riemann zeros.  Then:
\begin{enumerate}[label=(\alph*)]
  \item \textbf{(Criterion)} \emph{err-collapse equivalence:}
    $\mathrm{err}_\lambda \to 0$ if and only if conditions (i)$+$(ii)
    of Connes \cite[§6.6]{Connes2026} hold.  (Proved via the triangle
    inequality for Fourier-analytic distance.)
  \item \textbf{(Rigorous at every finite $N$)} \emph{Condition~(i):}
    real zeros of $\hat{\xi}_\lambda$ hold by the Connes--van Suijlekom
    Theorem~5.6 \cite{ConnesvS2025} at every finite~$N$, as long as the
    even/odd eigenvalue split is non-zero.
  \item \textbf{(Empirical)} At dps$\,= 250$--$350$, $\mathrm{err}_\lambda
    \sim 10^{-25\lambda}$ for the first 15 Riemann zeros, $\lambda=3,\ldots,9$.
\end{enumerate}
The remaining analytic open step is therefore identical to (ii-a): a
rigorous a-priori strip bound on $\hat{\xi}_\lambda$ from $QW_N$-coercivity,
as isolated in \cref{prop:iic-absorbed}.
\end{remark}

% =============================================================================
\section{Conclusion: The Zookeeper's Legacy}
\label{sec:conclusion}
% =============================================================================

Connes' spectral programme, culminating in the CCM Fourier basis
\cite{CCM2025}, provides the first computationally accessible
framework in which the low-energy mechanism behind the zeta zeros can be
tested mode by mode and eigenvalue by eigenvalue with arbitrary
precision.

Across the tested diagnostics, the framework shows a consistent
numerical pattern, with margins that range from comfortable to enormous.
The $15$-digit cancellation at Riemann zeros, the resolvent identity
$\alpha = R$, the Poisson kernel's near-perfect cluster localisation,
the sign coherence of the spectral multiplier, the $\eta \approx 1$
distribution---all of these are consequences of a single structural
fact: the rank-one pole term and the arithmetic perturbation are locked
in a rigid algebraic embrace.

\medskip
\noindent\textbf{What this paper proves analytically or conditionally reduces:}
\begin{enumerate}[label=(\roman*)]
\item The \emph{resolvent identity} (Theorem~\ref{thm:resolvent}):
  Riemann zeros are resonances $\alpha(\gamma) = R(\gamma, \wmin)$.
\item The \emph{$\eta$-parametrisation} (Theorem~\ref{thm:eta-equiv}):
  the local modewise inequality is equivalent to $\eta_j \in (0,2)$ for
  all regular modes.
\item The \emph{rank-one sign forcing} behind $\eta_j > 0$
  (Theorem~\ref{thm:eta-positive}), via the secular equation and Cauchy
  interlacing, with the finite multi-rank correction left to the
  conditional estimates.
\item The \emph{conditional three-statement microcluster closure}
  (Theorem~\ref{thm:MS2-closure}), which turns scalar secular
  cancellation, projected Poisson control, and the outer spectral gap
  into an explicit MS2 bound, assuming the three quantitative closure
  statements.
\item The \emph{Conditional reduction from microcluster closure to RH}
  (Theorem~\ref{thm:RH-conditional}), obtained by combining the
  conditional microcluster closure with the CCM transfer assumptions;
  this is a reduction, not a proof of RH.
\end{enumerate}

\noindent\textbf{What this paper verifies numerically:}
\begin{enumerate}[label=(\roman*)]
\item The upper bound $\eta_j < 2$ with margin $\geq 0.99$ (all 117+ modes).
\item The Abel--Stieltjes Tail-vs-Gap bound $< 1$ (all 117 modes).
\item The endgame constants $s_\lambda$, $p_\lambda$, and $g_*$ for
  $\lambda = 3, \ldots, 100$ ($N$ up to~$120$) in the tested range,
  with the closure bound 75--84\% tight.
\end{enumerate}

\noindent\textbf{What remains external to the internal closure argument:}
\begin{enumerate}[label=(O\arabic*)]
\item \textbf{Uniform C2ca constants:} the algebraic identities behind
  $s_\lambda$, $p_\lambda$, and $g_*$ are established or numerically very
  stable; the uniform $\lambda$-limit needed for MS2 remains a quantitative
  closure assumption.
\item \textbf{Truncation faithfulness and limit passage:} the present
  paper works inside the finite-dimensional CCM Fourier model.  To pass
  from the model to the full Weil positivity statement one still needs
  the external CCM transfer from the Fourier-Galerkin regime to the
  $\lambda\to\infty$ limit.
\item \textbf{MS1 / simple-even input:} the simplicity and even-sector
  properties used in the broader CCM programme are not reproved here.
\end{enumerate}

\bigskip

% =============================================================================
\section{Update (v1.6): Instrument Status and Contract-Form Applications,
June--August 2026}
\label{sec:update-v16}
% =============================================================================

This section is additive.  It records the state of the diagnostic
instrumentation surrounding the argument of this paper as of
August~2026, and it changes none of the statements, proofs, or numerical
claims above.  In particular, the conditional character of the reduction
in \cref{sec:RH} is unchanged, and nothing below is progress on the
missing step MS2 itself.

\subsection{Consistency instrumentation}

Since v1.5 the programme has continued to maintain a cross-route
consistency register: a shadow registry that records, for every route
under investigation, which statements are actually claimed, which are
merely diagnostic, and where two routes would contradict one another if
both were taken at face value.  Its purpose is bookkeeping and
falsification pressure, not derivation.  The registry produces no
theorem, and it is not evidence for any statement of this paper.

The register of the present programme stands at revision~V06.  The
sister programme on the abc conjecture maintains a separate register of
the same kind, whose canonical branch is v118 alongside an active
continuation branch v119.  The two registers are deliberately kept
apart; a merged cross-programme view is envisaged but has not been
carried out.

\subsection{Contract forms as sister instruments}

A second line of instrumentation was developed in that sister
programme, where a published chain of arguments had to be audited link
by link.  Three of its check forms are relevant here as \emph{forms},
independently of their arithmetic content:
\begin{enumerate}[label=(C\arabic*)]
\item a \emph{bridge certificate}, which fixes in advance what a claimed
  passage between two frameworks must deliver before it may be used
  \cite{Geiger2026bridge};
\item a \emph{ledger audit}, which applies such a certificate to a
  published chain and records, entry by entry, which links carry and
  which do not \cite{Geiger2026ledger};
\item a \emph{rigidity and effective-weight test}, which asks whether
  repairing a broken link is possible at all within the admissible
  class, or whether a rigidity statement excludes it
  \cite{Geiger2026rigidity}.
\end{enumerate}

Their relevance to the present paper is structural rather than
arithmetic.  The external inputs listed as (O1)--(O3) in
\cref{sec:conclusion} --- the uniform-in-$\lambda$ constants, the
truncation-faithfulness implication, and the passage from finite
approximants to the limit --- are precisely of the type these forms are
built to audit: transitions between a finite model and a limit object
which the source literature states as proof strategies rather than as
completed theorems.  Applying the certificate form to the spectral
transport step of this programme is therefore a natural next
instrument.  \emph{Status: envisaged, not carried out.}  No such audit
has been performed for the arguments of this paper, and no statement in
the present version depends on one.  The companion cartography that
situates the route of this paper among the alternatives has been
updated in parallel \cite{Geiger2026atlas}.

\subsection{Instrument status}

The following records the auxiliary spectral instrumentation examined
between June and August~2026.  Each line names what has been settled
\emph{and} what remains open; none of them is an input to any theorem
of this paper.
\begin{enumerate}[label=(I\arabic*)]
\item \textbf{Limit topology.}  For the ladder of spectral projectors the
  norm route is not merely unproved but \emph{excluded}: projectors of
  different rank remain at norm distance at least one, and the rank does
  increase along the ladder.  What remains admissible is strong
  convergence in the ambient space, with a possibly infinite limit rank.
  That topology is now named precisely; it has not been verified.
\item \textbf{Contour admissibility.}  From the quantities the instrument
  currently records, only an \emph{upper} bound on the distance between
  the contour and the spectrum can be derived.  A lower bound --- the
  direction actually required for a contour uniformly separated from the
  spectrum --- does not follow from them.  The question is therefore
  neither settled nor refuted, and further quantities must be measured
  before it can be decided either way.
\item \textbf{Window and seed.}  The selection of the spectral window and
  of its seed was made independent of the enumeration index, and the
  resulting finite-dimensional contract holds on the tested ladder.
  Uniformity in the ladder parameter and the behaviour in the limit are
  untouched by this and remain open.
\item \textbf{Selection family.}  The family produced by the selection
  rule is not canonical: different admissible rules yield different
  families.  The finite gate is passed; the limit gate is open, and at
  present no argument singles out one family as the right one.
\item \textbf{Inclusion angles.}  The angles the instrument reports are
  one-sided inclusion diagnostics, not symmetric measures of projector
  stability.  Angles between successive rungs do not decide the limit
  question --- under strong convergence they need not decrease, since
  new directions enter at every rung --- so their behaviour is evidence
  in neither direction.  The quantity that would be informative, taken
  against a fixed anchor space, has not yet been computed.
\end{enumerate}

\subsection{Status of this update}

Everything recorded in this section is instrument and validation
progress.  It is not a step toward MS2, it strengthens no statement of
\cref{sec:endgame} or \cref{sec:RH}, and it leaves the conditional
reduction exactly as conditional as before.  No closure statement has
become uniform, no external CCM input has been supplied, and no claim
about the Riemann Hypothesis is made or implied here.  The active
analytic fronts of the programme are pursued separately and will be
reported in a dedicated follow-up paper.

\bigskip

% =============================================================================
% ACKNOWLEDGMENTS
% =============================================================================
\subsection*{Acknowledgments}

The author thanks Alain Connes, Caterina Consani, and Henri Moscovici
for creating the framework that made this investigation possible.
The numerical computations were performed using \texttt{mpmath}
\cite{mpmath} at 50-digit precision.

\subsection*{Data availability}

All computational scripts and numerical logs are available at
\url{https://github.com/research-line/functional-stability-theory/tree/main/masters/zookeeper}.

\subsection*{AI disclosure}

AI tools (Anthropic Claude and OpenAI ChatGPT) were used for numerical
computation checks, algebraic verification support, proof-strategy
exploration, and manuscript preparation.  The tool-assisted steps are
documented in the supplementary materials.
The author retained final responsibility for the mathematical content and
the final editorial form.

% =============================================================================
% REFERENCES
% =============================================================================

\begin{thebibliography}{99}

\bibitem{CCM2025}
A.~Connes, C.~Consani, and H.~Moscovici,
\emph{Zeta spectral triples},
arXiv:2511.22755 (2025).

\bibitem{Connes1999}
A.~Connes,
\emph{Trace formula in noncommutative geometry and the zeros of the
Riemann zeta function},
Selecta Math.~(N.S.) \textbf{5} (1999), 29--106.

\bibitem{Weil1952}
A.~Weil,
\emph{Sur les ``formules explicites'' de la th{\'e}orie des nombres
premiers},
Meddelanden fr{\aa}n Lunds Univ.\ Mat.\ Sem., Tome suppl{\'e}mentaire
(1952), 252--265.

\bibitem{Li1997}
X.-J.~Li,
\emph{The positivity of a sequence of numbers and the Riemann hypothesis},
J.~Number Theory \textbf{65} (1997), 325--333.

\bibitem{Montgomery1973}
H.~L.~Montgomery,
\emph{The pair correlation of zeros of the zeta function},
Proc.\ Sympos.\ Pure Math.\ \textbf{24} (1973), 181--193.

\bibitem{Odlyzko1987}
A.~M.~Odlyzko,
\emph{On the distribution of spacings between zeros of the zeta function},
Math.\ Comp.\ \textbf{48} (1987), 273--308.

\bibitem{BerryKeating1999}
M.~V.~Berry and J.~P.~Keating,
\emph{The Riemann zeros and eigenvalue asymptotics},
SIAM Rev.\ \textbf{41} (1999), 236--266.

\bibitem{Connes2026}
A.~Connes,
\emph{The Riemann Hypothesis: Past, Present and a Letter Through Time},
arXiv:2602.04022 (2026).

\bibitem{ConnesvS2025}
A.~Connes and W.~D.~van Suijlekom,
\emph{Quadratic Forms, Real Zeros and Echoes of the Spectral Action},
Communications in Mathematical Physics \textbf{406} (2025), Article~312,
doi:10.1007/s00220-025-05493-1.

\bibitem{GolubVanLoan2013}
G.~H.~Golub and C.~F.~Van Loan,
\emph{Matrix Computations},
4th edition, Johns Hopkins University Press, 2013.

\bibitem{Geiger2026evendom}
L.~Geiger,
\emph{A Conditional Reduction of the Riemann Hypothesis to Even Dominance of the
Weil Quadratic Form},
Zenodo (2026), Concept DOI:~10.5281/zenodo.19764771.

\bibitem{mpmath}
F.~Johansson et al.,
\texttt{mpmath}: a Python library for arbitrary-precision floating-point
arithmetic (version 1.3),
\url{https://mpmath.org/} (2023).

\bibitem{Geiger2026bridge}
L.~Geiger,
\emph{The Bridge Certificate},
Zenodo (2026), Concept DOI:~10.5281/zenodo.21925803.

\bibitem{Geiger2026ledger}
L.~Geiger,
\emph{The Bridge Certificate Applied to the Joshi ATS Series:
A Ledger Audit},
Zenodo (2026), Concept DOI:~10.5281/zenodo.21951155.

\bibitem{Geiger2026rigidity}
L.~Geiger,
\emph{Repairing the ATS Chain?  A Rigidity Theorem for Product-Formula
Weights and an Effective-Weight Test for the Height Subchain},
Zenodo (2026), Concept DOI:~10.5281/zenodo.21952486.

\bibitem{Geiger2026atlas}
L.~Geiger,
\emph{From Landscape to Atlas: Multi-Route Cartography of an Ongoing
Expedition Toward the Riemann Hypothesis},
version~3.2, Zenodo (2026), DOI:~10.5281/zenodo.21952620.

\end{thebibliography}

\end{document}
